\chapter{Build a Connectivity Plan}\label{ch:nw:build-a-connectivity-plan}

A new firm will trade three venues: equities on an exchange in New Jersey, futures on an exchange outside Chicago, and a crypto
venue whose matching engine runs in a cloud region in Tokyo. Before it signs anything it has to say what latency it will get at
each, what it must buy, from whom, on what terms, and what it will cost a year. The book's last chapter answers with one
configuration and the book's own components: USD 1.02 million a year all in, three of four latency targets met, and a missed
one that a \omterm{def:nw:private-connectivity-to-crypto-venues:endpoint}{private endpoint} fixes for USD 123 a year.

Every earlier chapter supplied a piece. Chapter~10 gave the sites and the \omterm{def:nw:the-north-american-map:floor}{latency floors}; chapter~13 fibre routes and chapter~14
radio; chapters~2 to~8 the cage, cards, software path and servers; chapter~9 \omterm{def:nw:colocation-products-and-how-they-are-sold:colo}{colocation} prices; chapter~15 circuits and service
levels; chapters~16 to~18 the cloud; chapter~27 the vendor map; chapter~28 recovery. This chapter joins them in
\texttt{firm.connplan}, the book's capstone, and runs it on a plan whose prices are dated rows from the venues' own schedules.

\section{Requirements: strategies, venues and latency targets}

\begin{definition}[Connectivity plan, bill of materials]\label{def:nw:build-a-connectivity-plan:plan}
A \emph{connectivity plan}\index{connectivity plan} is a firm's document, kept as data, that states for each venue it trades the
latency it needs and how it will get it: the sites, routes, connections and services it will buy, their service levels and
recovery arrangements, and their cost. Its \emph{bill of materials}\index{bill of materials} is the itemised list of everything to
be bought, with quantity, supplier, price and the dated source of each price.
\end{definition}

The firm's strategies set the targets. Its equity strategy reacts to the equity venue's own book, so it colocates in Mahwah and
needs a wire-to-wire time under \qty{40}{\micro\second} at the 99th percentile. Its futures strategy hedges on CME: it needs Aurora's
prices in Mahwah within \qty{4.2}{\milli\second} and its orders in Aurora within \qty{8}{\milli\second}. Its crypto strategy trades a
venue in AWS's Tokyo region from instances there, and needs round trips to the venue under one millisecond at the 99th percentile
(all targets are the firm's, and assumptions here).

\omcode{code/networks/29-build-a-connectivity-plan/python/nw_plan.py}{17}{36}{The plan as data: venues with access and targets, and the footprint of priced lines.}\label{lst:nw:build-a-connectivity-plan:plan}

\section{The latency table}

Each access kind takes its latency from a component: a colocated path from chapter~2's cage and chapter~5's software path
(\texttt{firm.cagenet}, \texttt{firm.wirepath}, Book~13's measured stages); a radio route from its published figure; a fibre route
from the geodesic and a \omterm{def:nw:the-north-american-map:factor}{route factor} (\texttt{firm.fibreroute}); a cloud path from chapter~18's private paths
(\texttt{firm.privlink}).

\omcode{code/firm/connplan/firm_connplan.py}{100}{111}{One latency per venue, from the component that models its access.}\label{lst:nw:build-a-connectivity-plan:latency}

\begin{table}[ht]
\centering\small
\begin{tabular}{@{}llrrl@{}}
\toprule
Venue & Measure & Achieved & Target & \\
\midrule
NYSE equities (Mahwah) & wire-to-wire, 99th pct. & \qty{35.6}{\micro\second} & \qty{40}{\micro\second} & met \\
CME data, Aurora to Mahwah & one way (radio) & \qty{3.986}{\milli\second} & \qty{4.2}{\milli\second} & met \\
CME orders, Mahwah to Aurora & one way (fibre at 1.3) & \qty{7.48}{\milli\second} & \qty{8}{\milli\second} & met \\
Binance (AWS Tokyo) & round trip, 99th pct. & \qty{2.40}{\milli\second} & \qty{1}{\milli\second} & missed \\
\bottomrule
\end{tabular}
\caption{The plan's latency table. The radio figure is the fastest published Aurora--Mahwah route (chapter~10); NYSE's fee schedule
gives no latency for its wireless CME data. The fibre route's factor and the cloud path's edge time are assumptions. Data:
\texttt{nw\_plan.report()}.}\label{tab:nw:build-a-connectivity-plan:latency}
\end{table}

\begin{omfigure}
\begin{tikzpicture}
\begin{axis}[ybar,width=10cm,height=5.4cm,bar width=16pt,ymin=0,ymax=280,xtick=data,
  xticklabels from table={figdata/networks/29-build-a-connectivity-plan/latency.csv}{venue},
  ylabel={achieved latency, \% of target},nodes near coords,
  every node near coord/.append style={font=\scriptsize,anchor=north,yshift=-2pt},
  nodes near coords={\pgfmathprintnumber[fixed,fixed zerofill,precision=1]{\pgfplotspointmeta}},
  tick label style={font=\footnotesize},label style={font=\small},grid=major,grid style={gray!25},table/col sep=comma]
\draw[omThm,thick,dashed] (axis cs:-0.5,100) -- (axis cs:3.5,100);
\addplot[fill=omDef!55,draw=omDef] table[x=x,y=share]{figdata/networks/29-build-a-connectivity-plan/latency.csv};
\end{axis}
\end{tikzpicture}

\omcaption{Each venue's achieved latency as a share of its target; the dashed line is the target. Three venues are within it; the
Tokyo crypto venue, reached through the public internet edge, is at 240\%. Data: \texttt{fig\_plan.py},
\texttt{nw\_plan.report()}.}\label{fig:nw:build-a-connectivity-plan:latency}
\end{omfigure}

The missed target is the one the plan can fix cheaply. The configuration gives the crypto venue chapter~18's three paths (an
assumption about what it offers): its public endpoint through an \omterm{def:nw:where-the-crypto-venues-live:as}{edge network}, and a \omterm{def:nw:private-connectivity-to-crypto-venues:endpoint}{private endpoint} in the firm's zone or in another zone. The plan tries each and ranks those that meet the target by
cost.

\omcode{code/firm/connplan/firm_connplan.py}{154}{169}{The cheapest change that meets a missed cloud target: each offered path, tried and priced.}\label{lst:nw:build-a-connectivity-plan:fix}

A \omterm{def:nw:private-connectivity-to-crypto-venues:endpoint}{private endpoint} in the instance's zone brings the 99th-percentile round trip from \qty{2.40}{\milli\second} to
\qty{0.44}{\milli\second}, for USD 122.64 a year of endpoint hours (at chapter~18's price; its traffic charge adds to that with volume);
one in another zone meets the target too, at \qty{0.80}{\milli\second}, for the same fixed price plus cross-zone transfer. The fix is to
align the endpoint with the instance's zone identifier, chapter~16's lesson.

\section{The bill of materials}

\begin{dated}{2026-09}{Prices in the plan, from the venues' and the cloud's published schedules}\label{dat:nw:build-a-connectivity-plan:prices}
NYSE (schedule of 14 September 2026): dedicated cabinet 5\,000 USD initial and 1\,200 USD a kW a month at 4--8 kW; 10 Gb LCN and NMS
network connection 15\,000 USD initial and 24\,000 a month (22\,000 in NYSE's October 2025 schedule, as Nasdaq's filing cited it); cross
connect 500 and 600 a month; wireless CME data 5\,000 and 6\,000 a month; CME data feed over the IP network 3\,000 a month. Nasdaq's
comparable 10Gb Ultra connection: 18\,500 USD a month from January 2026. AWS Tokyo (price list of 25 September 2026): c7i.4xlarge
0.8988 USD an hour.
\end{dated}

\begin{table}[ht]
\centering\footnotesize
\begin{tabular}{@{}lrrrl@{}}
\toprule
Item & Qty & Monthly (USD) & Annual (USD) & Source \\
\midrule
NYSE dedicated cabinet (initial fee) & 1 & -- & 1\,667 & 29:F1 \\
NYSE cabinet power, 4--8 kW & 8 kW & 9\,600 & 115\,200 & 29:F1 \\
NYSE fibre cross connects & 4 & 2\,400 & 29\,467 & 29:F1 \\
NYSE 10 Gb LCN and NMS connection & 1 & 24\,000 & 293\,000 & 29:F1 \\
NYSE IP network 1 Gb (backup) & 1 & 2\,500 & 30\,833 & 15:F4 \\
Wavelength Mahwah--Aurora (CME orders) & 1 & 20\,000 & 243\,333 & assumed \\
NYSE wireless CME data & 1 & 6\,000 & 73\,667 & 29:F1 \\
CME data feed over the IP network & 1 & 3\,000 & 36\,000 & 29:F1 \\
AWS Tokyo c7i.4xlarge (730 h) & 2 & 1\,312 & 15\,747 & 29:F2 \\
\bottomrule
\end{tabular}
\caption{The plan's \omterm{def:nw:build-a-connectivity-plan:plan}{bill of materials}; one-time fees amortised over 36 months in the annual column. The wavelength has no published
price: the plan carries an assumed one, and its own check flags it. Data: \texttt{code/firm/connplan/data/cost\_table.csv}.}\label{tab:nw:build-a-connectivity-plan:bom}
\end{table}

Two checks run before the plan is trusted: every price row must name its ledger source, and none may be older than 60 days. On the day
of writing, one row fails the first check, the Mahwah--Aurora wavelength, whose USD 20\,000 a month is an assumption the firm must replace
with a carrier's quote; none fails the second. The table is also the plan's export: \texttt{export\_cost\_table} writes it with its
sources and dates for the firm's budget owners (Book~16 reads it).

\section{Contracts and service levels}

The plan's access to the equity venue is one 10 Gb connection, with a 1 Gb IP-network circuit as a backup. With each circuit failing
once in 4\,000 hours and repaired in four (assumptions), one connection alone is available 99.900\% of the time; with the backup and a
common-mode event (a shared duct or room) twice in ten years for eight hours, 99.982\% (\texttt{firm.slamodel}, chapter~15). A second
10 Gb connection beside the backup would add USD 366\,250 a year to the budget with the recovery share: the price of keeping full
bandwidth through a failure. The fee schedule read for the plan mentions no
service-level credits for these circuits; they belong in the contract.

\section{The annual budget and its sensitivities}

\begin{omfigure}
\begin{tikzpicture}
\begin{axis}[xbar,width=10cm,height=5cm,bar width=9pt,xmin=0,xmax=650,ytick=data,y dir=reverse,
  yticklabels from table={figdata/networks/29-build-a-connectivity-plan/budget.csv}{category},
  xlabel={thousand USD a year},nodes near coords,every node near coord/.append style={font=\scriptsize},
  nodes near coords={\pgfmathprintnumber[fixed,fixed zerofill,precision=1]{\pgfplotspointmeta}},
  tick label style={font=\footnotesize},label style={font=\small},grid=major,grid style={gray!25},table/col sep=comma]
\addplot[fill=omThm!45,draw=omThm] table[y=x,x=kusd]{figdata/networks/29-build-a-connectivity-plan/budget.csv};
\end{axis}
\end{tikzpicture}

\omcaption{The plan's annual budget by category, USD 1.02 million in all. Connectivity dominates; disaster recovery is a quarter of the
on-site \omterm{def:nw:colocation-products-and-how-they-are-sold:colo}{colocation} and connectivity (chapter~28's warm design, an assumption). Data: \texttt{fig\_plan.py},
\texttt{nw\_plan.report()}.}\label{fig:nw:build-a-connectivity-plan:budget}
\end{omfigure}

The plan costs USD 68\,812 a month plus 39\,500 in one-time fees: USD 1\,017\,289 a year with the fees amortised over three years and the
recovery site at a quarter of the on-site spend. Connectivity is 56\% of it, and one connection, the 10 Gb LCN and NMS link, is 29\%
alone. The sensitivities show where the budget moves: a second LCN connection adds USD 366\,250 a year; four more kW of power 72\,000; a
wavelength priced 50\% above the assumption 152\,083; doubling the Tokyo instances only 15\,747; dropping the backup circuit saves 38\,542
and costs the availability above. The cloud leg, the one that missed its target, is the cheapest to fix and to grow.

\begin{method}[Building and keeping a connectivity plan]\label{met:nw:build-a-connectivity-plan:method}
\begin{enumerate}
\item Write the requirements as data: each venue, what each strategy needs from it, and the measure (percentile, one way or round trip).
\item Compute each latency with the component that models its access; compare with the target and flag every miss.
\item For each miss, try the offered alternatives and rank those that meet the target by cost.
\item Price every line from a dated, sourced row; flag unsourced and stale rows and replace assumptions with quotes before signing.
\item Add service levels and the recovery design; compute the budget and its sensitivities; export the cost table.
\item Re-run the plan at every price change, venue move or new strategy: the plan is a program, not a document.
\end{enumerate}
\end{method}

\section{Tutorial: three venues, one budget}\label{tut:nw:build-a-connectivity-plan:tut}

\begin{tutorial}
\textbf{Goal.} Run the whole plan from one configuration and change one requirement.
\textbf{End state:} \cref{tab:nw:build-a-connectivity-plan:latency,tab:nw:build-a-connectivity-plan:bom} and
\cref{fig:nw:build-a-connectivity-plan:budget}.
\begin{tutsteps}
\item \textbf{Configuration.} \texttt{nw\_plan.VENUES} and \texttt{LINES} (\cref{lst:nw:build-a-connectivity-plan:plan}) with
  \texttt{firm\_connplan.load\_prices}.
\item \textbf{Latency.} \texttt{latency\_table} (\cref{lst:nw:build-a-connectivity-plan:latency}); \texttt{cheapest\_fixes}
  (\cref{lst:nw:build-a-connectivity-plan:fix}).
\item \textbf{Money.} \texttt{bom}, \texttt{budget}, \texttt{check\_prices}, \texttt{nw\_plan.sensitivities} and
  \texttt{export\_cost\_table}.
\end{tutsteps}
\textbf{What to change next.} Tighten the equity target to \qty{20}{\micro\second} and let the plan propose chapter~7's hardware path;
move the futures strategy to Aurora and compare the budgets.
\end{tutorial}

\section{Build: the connectivity plan}\label{bld:nw:build-a-connectivity-plan:connplan}

\begin{build}
\bfield{Purpose} The firm's \omterm{def:nw:build-a-connectivity-plan:plan}{connectivity plan} as a program: latency table, \omterm{def:nw:build-a-connectivity-plan:plan}{bill of materials}, contracts, recovery line, budget,
sensitivities, fixes, checks and the exported cost table. The book's capstone.
\bfield{Interface} \texttt{firm\_connplan}: \texttt{PriceRow}, \texttt{load\_prices}, \texttt{Venue}, \texttt{Line}, \texttt{Plan},
\texttt{latency\_us}, \texttt{latency\_table}, \texttt{bom}, \texttt{budget}, \texttt{check\_prices}, \texttt{cheapest\_fixes},
\texttt{availability}, \texttt{export\_cost\_table}; the exported table's columns in its \texttt{README.md}.
\bfield{Rules} Every price row names its ledger source and date, or is flagged; latencies come from the book's components; targets and
assumptions are stated in the configuration.
\bfield{Acceptance tests} \texttt{code/firm/connplan/tests/}: price rows, a budget by hand, the checks and the fix on a small plan, a fibre
route against its floor, and the exported table.
\bfield{Stretch} Every access kind for every venue, with its alternatives priced; currency conversion with dated rates; the hardware
path as an option for colocated venues.
\end{build}

\begin{omsources}
\item NYSE, Connectivity Fees and Charges (September 2026); Nasdaq, SR-NASDAQ-2025-089; AWS EC2 price list (Tokyo, September 2026).
\item This book's chapters~2 to~28 and their ledgers.
\end{omsources}

\section{Exercises}

\begin{exercise}[$\star$]\label{exo:nw:build-a-connectivity-plan:1}
What is a \omterm{def:nw:build-a-connectivity-plan:plan}{bill of materials}, and what must each of its rows carry in this book's plan?
\end{exercise}

\begin{exercise}[$\star$]\label{exo:nw:build-a-connectivity-plan:2}
What does an 8 kW dedicated cabinet cost a month in Mahwah on the September 2026 schedule?
\end{exercise}

\begin{exercise}[$\star$]\label{exo:nw:build-a-connectivity-plan:3}
Why does the plan flag the Mahwah--Aurora wavelength?
\end{exercise}

\begin{exercise}[$\star\star$]\label{exo:nw:build-a-connectivity-plan:4}
How much of the annual budget is the 10 Gb LCN and NMS connection, and how much did its monthly price rise between NYSE's October 2025 and
September 2026 schedules?
\end{exercise}

\begin{exercise}[$\star\star$]\label{exo:nw:build-a-connectivity-plan:5}
Why does a \omterm{def:nw:private-connectivity-to-crypto-venues:endpoint}{private endpoint} in another zone meet the Tokyo target but still cost more than one in the instance's zone?
\end{exercise}

\begin{exercise}[$\star\star$]\label{exo:nw:build-a-connectivity-plan:6}
Which single change moves the budget most, and why is it not obviously worth it?
\end{exercise}

\begin{exercise}[$\star\star\star$]\label{exo:nw:build-a-connectivity-plan:7}
\emph{Coding.} With \texttt{firm\_connplan}, what one-way latency does the CME order route reach with a \omterm{def:nw:the-north-american-map:factor}{route factor} of 1.2 instead of 1.3?
\end{exercise}

\begin{exercise}[$\star\star\star$]\label{exo:nw:build-a-connectivity-plan:8}
\emph{Find the flaw.} ``Our plan costs USD 1.02 million a year; we can sign today.''
\end{exercise}

\section{Problem: Three Venues, One Budget}

\begin{problem}[{Weekend problem --- the capstone plan}]\label{pb:nw:build-a-connectivity-plan:1}
The firm of the chapter must present its plan to its board. Use \texttt{firm.connplan} and the chapter's configuration.

\medskip\noindent\textbf{Part I --- Latency.}
\begin{enumerate}
\item What does each venue achieve against its target?
\item Which components computed each figure?
\item Which of the figures are published and which modelled?
\item What is missed, by how much, and why?
\end{enumerate}

\medskip\noindent\textbf{Part II --- The fix.}
\begin{enumerate}[resume]
\item Which options meet the missed target, and at what latency?
\item What does the cheapest cost a year?
\item What must the firm check in its cloud account for the fix to work?
\item Why is the fix so cheap compared with the rest of the plan?
\end{enumerate}

\medskip\noindent\textbf{Part III --- Money.}
\begin{enumerate}[resume]
\item What are the monthly, one-time and annual costs?
\item How is the annual cost split by category?
\item Which row fails the checks, and what should the firm do?
\item What do the sensitivities show?
\end{enumerate}

\medskip\noindent\textbf{Part IV --- The verdict.}
\begin{enumerate}[resume]
\item State the \emph{named result}: the annual all-in connectivity cost of entering the three venues, the latency achieved at each
  against its target, and the cheapest change that meets the one missed target.
\item What does the backup circuit buy?
\item What would moving the futures strategy to Aurora change?
\item What should the board ask before approving?
\item How often should the plan be re-run?
\item What does Book~16 take from the plan?
\item What part of the plan would you check first if a latency target were missed in production?
\item In one sentence: what is a \omterm{def:nw:build-a-connectivity-plan:plan}{connectivity plan} for?
\end{enumerate}
\end{problem}

\section{Interview questions}

\begin{interviewq}[$\star$ \iqroles{developer}]\label{iq:nw:build-a-connectivity-plan:1}
What goes into a trading firm's \omterm{def:nw:build-a-connectivity-plan:plan}{connectivity plan}?
\end{interviewq}

\begin{interviewq}[$\star\star$ \iqroles{developer}]\label{iq:nw:build-a-connectivity-plan:2}
How would you estimate the latency from New Jersey to a futures exchange near Chicago before buying anything?
\end{interviewq}

\begin{interviewq}[$\star\star$ \iqroles{developer, researcher}]\label{iq:nw:build-a-connectivity-plan:3}
Your plan misses one latency target. How do you decide what to change?
\end{interviewq}

\begin{interviewq}[$\star\star$ \iqroles{developer}]\label{iq:nw:build-a-connectivity-plan:4}
How do you keep a plan's prices trustworthy over time?
\end{interviewq}

\begin{interviewq}[$\star\star$ \iqroles{researcher}]\label{iq:nw:build-a-connectivity-plan:5}
The board asks whether the second 10 Gb connection is worth USD 366\,000 a year. How do you answer?
\end{interviewq}

\begin{interviewq}[$\star\star\star$ \iqroles{developer, researcher}]\label{iq:nw:build-a-connectivity-plan:6}
Design the connectivity of a new firm trading U.S. equities, CME futures and a crypto venue in Tokyo, from requirements to budget.
\end{interviewq}
